% Generated from ../chapters/24-how-this-framework-could-be-wrong.md
\chapter{How This Framework Could Be Wrong}

A framework becomes scientific when it risks failure.

\section{Error 1: ordinary adaptation mistaken for deep reorganization}

The reported changes may consist mainly of posture, muscle tone, body composition, attention, and pain modulation. These are valuable but do not imply broad structural restoration.

\section{Error 2: selection and memory bias}

Unusual improvements may be remembered while plateaus and regressions disappear from the narrative. Long duration can increase confidence without increasing measurement quality.

\section{Error 3: sounds mistaken for anatomical change}

Joint and tendon sounds may be benign mechanical events unrelated to remodeling or healing.

\section{Error 4: healthy attractors are weaker than proposed}

Adult tissues may not retain enough developmental flexibility for meaningful convergence after injury, aging, or structural loss.

\section{Error 5: positive feedback saturates early}

Better function may improve future adaptation only within a limited range. The system may plateau far before endogenous repair exceeds aging.

\section{Error 6: tissue-specific barriers dominate}

Some systems may reorganize extensively while others remain limited by missing cells, scar, mutations, crosslinks, or lost connectivity.

\section{Error 7: practices work through simpler mechanisms}

Benefits may be fully explained by exercise, stress reduction, sleep, expectation, social factors, and ordinary rehabilitation without requiring a new systemic theory.

This would not make the practices useless. It would make the stronger interpretation unnecessary.

\section{Error 8: the framework causes harm}

A person may reject treatment, fast excessively, overtrain, interpret neurological symptoms as healing, or force joints to release. Any framework that encourages such behaviour has failed as practice even if parts of its theory are correct.

\section{Falsifying observations}

The strongest version would be weakened if:

\begin{itemize}
\tightlist
\item
  long-term practitioners do not show greater structural or physiological plasticity than matched active controls;
\item
  measured changes remain within ordinary variability;
\item
  improvement does not increase future adaptive capacity;
\item
  cycles do not outperform simpler training and recovery;
\item
  claimed tissue changes cannot be independently verified;
\item
  benefits disappear when expectancy and attention are controlled.
\end{itemize}

The purpose of listing these possibilities is not defensive modesty. It is to make progress possible.

\begin{center}\rule{0.5\linewidth}{0.5pt}\end{center}
